\chapter{Continuation Algorithms}

The committed 3D benchmark follows one continuation path only:
\[
\text{SSR} \rightarrow \text{indirect method} \rightarrow \text{classic mode}
\rightarrow \text{secant predictor} \rightarrow \text{legacy omega controller}.
\]

This chapter now follows only that path.

\section{Mainline Continuation Contract}

The benchmark defines or inherits the following effective continuation settings:
\begin{itemize}
\item \code{analysis = ssr},
\item \code{method = indirect},
\item \code{predictor = secant},
\item \code{omega\_step\_controller = legacy},
\item \code{continuation\_mode = classic},
\item \code{continuation\_secant\_correction\_mode = none},
\item \code{continuation\_first\_newton\_warm\_start\_mode = none},
\item \code{step\_length\_cap\_mode = none},
\item \code{fine\_switch\_mode = none}.
\end{itemize}

That list is the main reason the continuation chapter is cleaner than before:
many repository features exist, but the default benchmark uses a fairly narrow
subset of them.

\section{Mainline Call Chain}

The active continuation chain is:
\begin{enumerate}
\item \code{SSR\_indirect\_continuation(...)} in
\petscmod{continuation/indirect.py},
\item initialization of the first accepted points,
\item a secant predictor for the next candidate state,
\item a nested Newton solve through \code{newton\_ind\_ssr(...)},
\item acceptance or rejection of the candidate point,
\item update of continuation history and the solver deflation basis.
\end{enumerate}

This is the modern PETSc equivalent of MATLAB
\matlabpkg{+CONTINUATION/SSR\_indirect\_continuation.m}.

\section{Accepted-Step Lifecycle}

One accepted continuation step on the current path is:
\begin{enumerate}
\item choose the next target \(\omega\),
\item build the secant predictor \((U_{\mathrm{ini}}, \lambda_{\mathrm{ini}})\),
\item solve the nested Newton problem for that target,
\item accept the point if Newton succeeds,
\item append the accepted state to the branch history,
\item update the solver recycle basis.
\end{enumerate}

Rejected attempts still exist, but the basic branch-tracking logic is the same
one MATLAB authors already know.

\section{What The Legacy Controller Means Here}

Because the benchmark keeps \code{omega\_step\_controller = legacy}, the current
default path uses the older, simpler continuation-step adaptation logic rather
than the newer adaptive controller family. In practical terms:
\begin{itemize}
\item step growth and shrinkage are still driven by continuation success and
Newton effort,
\item there is no separate adaptive efficiency controller active,
\item there is no fine-switch phase transition active.
\end{itemize}

\section{Mainline History Output}

The current path still records the accepted branch history familiar from MATLAB:
\begin{itemize}
\item \(\lambda\) history,
\item \(\omega\) history,
\item displacement history,
\item accepted-step and attempt statistics,
\item linear and nonlinear work counters.
\end{itemize}

The difference is that PETSc stores these in explicit structured outputs rather
than a mostly monolithic MATLAB struct.

\section{Where To Edit}

\begin{itemize}
\item \petscmod{continuation/indirect.py} for the mainline SSR driver,
\item \petscmod{nonlinear/newton.py} for the nested Newton solve it calls,
\item \petscmod{core/run\_config.py} if the mainline continuation controls
should be exposed differently through TOML.
\end{itemize}

\section{Where The Other Continuation Features Went}

Direct SSR, limit load, adaptive controllers, reduced predictors, secant
corrections, fine-switch logic, first-Newton warm starts, and streaming
microsteps are moved to Appendix~\ref{app:continuation}.
